\documentclass[border=10pt]{standalone}
\usepackage[no-math]{fontspec}
\usepackage[brazil]{babel}
\usepackage{amsmath,amssymb}
\IfFileExists{C:/Windows/Fonts/Lato_Regular.ttf}{%
  \setmainfont{Lato}[
    Path = C:/Windows/Fonts/,
    UprightFont = Lato_Regular.ttf,
    BoldFont = Lato_Bold.ttf,
    ItalicFont = Lato_Italic.ttf,
    BoldItalicFont = Lato_BoldItalic.ttf
  ]
  \setsansfont{Lato}[
    Path = C:/Windows/Fonts/,
    UprightFont = Lato_Regular.ttf,
    BoldFont = Lato_Bold.ttf,
    ItalicFont = Lato_Italic.ttf,
    BoldItalicFont = Lato_BoldItalic.ttf
  ]
}{%
  \RequirePackage[default]{lato}
}
\renewcommand{\familydefault}{\sfdefault}

\usepackage{xcolor}
\usepackage{tikz}
\usepackage{pgfplots}
\pgfplotsset{compat=1.18}
\usetikzlibrary{arrows.meta,positioning,shapes.geometric,fit,backgrounds,calc,decorations.pathreplacing,patterns}

% Paleta institucional PLANTERR / UEFS (alinhada ao estilo planterr-slides.css / 09a)
\definecolor{etNavy}{HTML}{1F4268}
\definecolor{etNavyDark}{HTML}{152E4A}
\definecolor{etGold}{HTML}{B38A3D}
\definecolor{etGoldDark}{HTML}{8C6927}
\definecolor{etTeal}{HTML}{0E7C7B}
\definecolor{etForest}{HTML}{2D6A4F}
\definecolor{etAlert}{HTML}{9E2A2B}
\definecolor{etInk}{HTML}{26313D}
\definecolor{etMuted}{HTML}{5B6573}
\definecolor{etBorder}{HTML}{D5DEE8}

\definecolor{bgNavy}{HTML}{EDF2F7}
\definecolor{bgGold}{HTML}{F8F3E9}
\definecolor{bgTeal}{HTML}{EBF6F5}
\definecolor{bgForest}{HTML}{EEF5F0}
\definecolor{bgAlert}{HTML}{FDF2F2}
\definecolor{bgCanvas}{HTML}{FFFFFF}

\tikzset{
  >=Stealth,
  etBox/.style={
    draw=#1, fill=white, rounded corners=5pt, line width=1.15pt,
    inner sep=7pt, align=center, font=\footnotesize\sffamily, text=etInk
  },
  etTint/.style={
    draw=#1, fill=#1!8!white, rounded corners=5pt, line width=1.15pt,
    inner sep=7pt, align=center, font=\footnotesize\sffamily, text=etInk
  },
  etBadge/.style={
    fill=#1, text=white, rounded corners=9pt, inner xsep=8pt, inner ysep=3pt,
    font=\scriptsize\bfseries\sffamily
  },
  etTitle/.style={font=\normalsize\bfseries\sffamily, text=etNavy, align=center},
  etSub/.style={font=\footnotesize\sffamily, text=etMuted, align=center},
  etArrow/.style={->, >=Stealth, line width=1.35pt},
  etDArrow/.style={<->, >=Stealth, line width=1.35pt}
}


\begin{document}
\begin{tikzpicture}[font=\sffamily]

% Título
\node[etTitle, font=\large\bfseries\sffamily, text=etNavy] (title) at (8.0, 5.05)
  {Inovação Aberta (\textit{Open Innovation}) e Trajetória de \textit{Spin-offs} Acadêmicas};

% Estilo local para cards em 4 colunas
\tikzset{
  oiCard/.style={
    draw=#1, fill=white, rounded corners=6pt, line width=1.2pt,
    minimum width=3.70cm, minimum height=2.95cm, inner sep=0pt, anchor=north
  },
  oiHead/.style={
    fill=#1, text=white, rounded corners=4pt,
    minimum width=3.44cm, minimum height=0.66cm,
    align=center, font=\scriptsize\bfseries\sffamily
  },
  oiBody/.style={
    text width=3.26cm, align=left, font=\scriptsize\sffamily, text=etInk,
    inner sep=2pt
  }
}

% Card 1
\node[oiCard=etNavy] (c1) at (1.95, 4.30) {};
\node[oiHead=etNavy, anchor=north] (h1) at ([yshift=-0.13cm]c1.north) {1. Pesquisa \& NIT};
\node[oiBody, anchor=north] at ([yshift=-0.22cm]h1.south) {\raggedright\hyphenpenalty=10000%
  \textbullet~Pesquisa universitária\\[5pt]
  \textbullet~Proteção de PI (INPI)\\[5pt]
  \textbullet~Marco Legal de CT\&I
};

% Card 2
\node[oiCard=etGold] (c2) at (5.98, 4.30) {};
\node[oiHead=etGold, anchor=north] (h2) at ([yshift=-0.13cm]c2.north) {2. Validação \& Fomento};
\node[oiBody, anchor=north] at ([yshift=-0.22cm]h2.south) {\raggedright\hyphenpenalty=10000%
  \textbullet~Prova de conceito (TRL)\\[5pt]
  \textbullet~Incubadora / Parque Tec.\\[5pt]
  \textbullet~Subvenção (FINEP/FAPs)
};

% Card 3
\node[oiCard=etTeal] (c3) at (10.02, 4.30) {};
\node[oiHead=etTeal, anchor=north] (h3) at ([yshift=-0.13cm]c3.north) {3. \textit{Spin-off} / Licenciamento};
\node[oiBody, anchor=north] at ([yshift=-0.22cm]h3.south) {\raggedright\hyphenpenalty=10000%
  \textbullet~Nova EBT ou transferência\\[5pt]
  \textbullet~Royalties ou participação (\textit{Equity})
};

% Card 4
\node[oiCard=etNavyDark] (c4) at (14.05, 4.30) {};
\node[oiHead=etNavyDark, anchor=north] (h4) at ([yshift=-0.13cm]c4.north) {4. Impacto Territorial};
\node[oiBody, anchor=north] at ([yshift=-0.22cm]h4.south) {\raggedright\hyphenpenalty=10000%
  \textbullet~Empregos qualificados\\[5pt]
  \textbullet~Adensamento de cadeia\\[5pt]
  \textbullet~Retenção de talentos
};

% Setas sequenciais
\draw[etArrow, etNavy] (c1.east) -- (c2.west);
\draw[etArrow, etGoldDark] (c2.east) -- (c3.west);
\draw[etArrow, etTeal] (c3.east) -- (c4.west);

% Banner inferior
\node[
  draw=etNavy, fill=bgNavy, rounded corners=6pt, line width=1.15pt,
  text width=15.05cm, inner xsep=14pt, inner ysep=9pt,
  align=left, font=\footnotesize\sffamily, text=etInk, anchor=north
] (banner) at (8.0, 0.92) {%
  \textcolor{etNavy}{\textbf{Open Innovation (Chesbrough, 2003):}} fronteiras permeáveis permitem que ideias internas gerem novas empresas (\textit{spin-offs}, \textit{inside-out}) e tecnologias externas entrem via parcerias (\textit{outside-in}).
};
% Barra lateral de destaque no banner
\fill[etNavy, rounded corners=2pt] ([xshift=1.5pt, yshift=-3pt]banner.north west) rectangle ([xshift=5.5pt, yshift=3pt]banner.south west);

\end{tikzpicture}
\end{document}
